% ReportKit four-page algorithm guide reference spread.
% This is an integration fixture for the 2026-09-18 algorithm-visuals fix
% sprint. It deliberately uses the repaired primitives together at guide
% scale; it is not a new public primitive or a replacement for the full guide.
\documentclass{reportkit}
\usepackage{reportkit-longform}

\setreportkitleftheader{REPORTKIT / ALGORITHM GUIDE}
\setreportkitfooter{Algorithm visuals reference spread}
\setreportkitversion{v1.3.0}
\setreportkitsubject{Algorithm visuals integration fixture}
\setreportkitkeywords{ReportKit, algorithm visuals, guide fixture}
\title{Algorithm visuals reference spread}
\author{ReportKit acceptance fixture}
\date{19 September 2026}

\begin{document}
\RKFrontMatterBegin

% The longform title helper intentionally ends with \clearpage. This fixture
% needs the compact navigation and grouped pattern table on the same opening
% page, so the small title-page composition stays local to this visual artefact
% while reusing the public compact-contents helper.
\thispagestyle{empty}
\begin{center}
  \vspace*{0.08\textheight}
  {\sffamily\fontsize{8.5}{10}\selectfont\bfseries\color{LinkBlue}
    REPORTKIT / ALGORITHM GUIDE\par}
  \vspace{7mm}
  {\sffamily\bfseries\fontsize{23}{27}\selectfont\color{Ink}
    Algorithm visuals\par}
  {\sffamily\bfseries\fontsize{23}{27}\selectfont\color{Ink}
    reference spread\par}
  \vspace{4mm}
  {\sffamily\color{Muted}Four pages of guide-scale state and execution views\par}
  \vspace{7mm}
  {\sffamily\fontsize{8}{10}\selectfont\color{Muted}
    \RKCompactContents{Two pointers and sliding window, Stack and queue, Topological sort}\par}
  \vspace{9mm}
  \begingroup
  \sffamily\fontsize{8.2}{10}\selectfont
  \renewcommand{\arraystretch}{1.35}
  \begin{tabular}{@{}p{0.28\textwidth}p{0.62\textwidth}@{}}
    \toprule
    \textbf{Pattern group} & \textbf{Question the figure should answer} \\
    \midrule
    Two pointers & Which boundary moved, and what remains active? \\
    Sliding window & Which item entered or left the window? \\
    Traversal & What is next on the queue or stack? \\
    Topological sort & Which dependency node can execute next? \\
    \bottomrule
  \end{tabular}
  \endgroup
  \vfill
  {\sffamily\fontsize{8}{10}\selectfont\color{Muted}
    ReportKit visual release-gate fixture · default theme\par}
\end{center}
\clearpage
\RKMainMatterBegin

\RKSectionOpener{Two pointers and sliding window}
The state view carries the movement semantics; the short code units stay with
their introducing prose through the explicit \texttt{keep=auto} contract.

\noindent
\begin{minipage}[t]{0.48\linewidth}
\begin{diagram}[type=array,width=\linewidth,caption={Two pointers bind the active range.},description={Two pointers example: the left pointer L is at index zero, the right pointer R is at index four, and the active range is the candidate interval.}]
\begin{arraystate}[indices=auto,legend=auto]
  \cell[state=current]{1}\cell[state=active]{3}\cell[state=active]{5}\cell[state=active]{7}\cell[state=current]{9}
  \pointer[role=left,below]{L}{0}
  \pointer[role=right,below]{R}{4}
  \range[state=active]{0}{4}
\end{arraystate}
\end{diagram}

\begin{codeblock}[Two-pointer][keep=auto]{Python}
def pair_sum(nums, target):
    left, right = 0, len(nums) - 1
    return nums[left], nums[right]
\end{codeblock}
\end{minipage}\hfill
\begin{minipage}[t]{0.48\linewidth}

\begin{diagram}[type=array,width=\linewidth,caption={A window records entry and exit.},description={Sliding window example: item seven enters at index five while duplicate item seven leaves at index two; left and right window roles show the inclusive active band.}]
\begin{windowstate}[legend=auto]
  \values{4,2,7,1,3,7}
  \window{3}{4}
  \entering{5}{7}
  \leaving{2}{7}
\end{windowstate}
\end{diagram}

\begin{codeblock}[Window update][keep=auto]{Python}
def advance(window, entering, leaving):
    window.add(entering)
    window.remove(leaving)
\end{codeblock}
\end{minipage}

\begin{diagram}[type=trace,width=.98\linewidth,source={}]
\begin{algorithmtrace}[mode=strip,columns=2,snapshot width=1.5,gap=1.1]
  \snapshot{Before}{%
    \begin{arraystate}[indices=false]
      \cell{1}\cell{3}\cell{5}
      \pointer[role=left,below]{L}{0}
    \end{arraystate}%
  }
  \tracetransition{advance L}
  \snapshot{After}{%
    \begin{arraystate}[indices=false]
      \cell{1}\cell{3}\cell{5}
      \pointer[role=left,below]{L}{1}
    \end{arraystate}%
  }
\end{algorithmtrace}
\end{diagram}

\clearpage
\RKSectionOpener{Stack and queue}
Stack and queue share one reading unit here: the container chrome shows the
direction of work, while the adjacent graph compositions show how traversal
state and worklists reinforce one another.

\begin{diagram}[type=stack,width=\textwidth,caption={LIFO and FIFO worklists at guide scale.},description={Stack and queue example: the stack top is the next depth-first item, and the queue front is the next breadth-first item.}]
\begin{stackstate}
  \push[state=visited]{A}
  \push[state=frontier]{B}
\end{stackstate}
\begin{scope}[xshift=5.2cm]
  \begin{queuestate}
    \enqueue[state=current]{A}
    \enqueue[state=frontier]{B}
  \end{queuestate}
\end{scope}
\end{diagram}

\begin{diagram}[type=graph,width=\textwidth,caption={BFS composes graphstate with queuestate.},description={BFS composition example: the graph's frontier node is the next candidate, and the queue holds that candidate after the current node.}]
\begin{graphstate}[columns=2,x spacing=2.3,y spacing=1.55]
  \graphnode[state=visited]{bfs-a}{Start}
  \graphnode[state=current]{bfs-b}{Current}
  \graphnode[state=frontier]{bfs-c}{Queued}
  \graphnode[state=unseen]{bfs-d}{Unseen}
  \graphedge{bfs-a}{bfs-b}
  \graphedge{bfs-b}{bfs-c}
  \graphedge{bfs-b}{bfs-d}
\end{graphstate}
\begin{scope}[xshift=7.0cm,yshift=-1.0cm]
  \begin{queuestate}
    \enqueue[state=current]{Current}
    \enqueue[state=frontier]{Queued}
  \end{queuestate}
\end{scope}
\end{diagram}

\begin{diagram}[type=graph,width=\textwidth,caption={DFS composes graphstate with stackstate.},description={DFS composition example: the graph's frontier node is pushed onto the call stack below the current node.}]
\begin{graphstate}[columns=2,x spacing=2.3,y spacing=1.55]
  \graphnode[state=visited]{dfs-a}{Start}
  \graphnode[state=current]{dfs-b}{Current}
  \graphnode[state=frontier]{dfs-c}{Pushed}
  \graphnode[state=unseen]{dfs-d}{Unseen}
  \graphedge{dfs-a}{dfs-b}
  \graphedge{dfs-b}{dfs-c}
  \graphedge{dfs-b}{dfs-d}
\end{graphstate}
\begin{scope}[xshift=7.0cm,yshift=-1.0cm]
  \begin{stackstate}
    \push[state=visited]{Current}
    \push[state=frontier]{Pushed}
  \end{stackstate}
\end{scope}
\end{diagram}

\clearpage
\RKSectionOpener{Topological sort}
The dependency layout makes prerequisite distance visible from left to right.
The ready queue is the execution decision: its zero-indegree badge agrees with
the node that can execute next.

\begin{diagram}[type=dag,width=\textwidth,caption={A pipeline ready queue under Kahn's algorithm.},description={Topological sort example: raw orders is resolved, clean orders is the next executable frontier node with zero indegree, and daily sales remains unseen with one dependency. The ready queue contains clean orders.}]
\begin{dagstate}[layout=dependency]
  \dagnode[indegree=0,state=resolved]{raw}{raw\_orders}
  \dagnode[indegree=0,state=frontier]{clean}{clean\_orders}
  \dagnode[indegree=1,state=unseen]{sales}{daily\_sales}
  \dependency{raw}{clean}
  \dependency{clean}{sales}
  \readyqueue{clean}
\end{dagstate}
\end{diagram}

\begin{codeblock}[Kahn's ready step][keep=auto]{Python}
ready = deque([clean_orders])
node = ready.popleft()
run(node)
\end{codeblock}

\end{document}
